\documentclass[11pt]{article}
\usepackage[margin=1in]{geometry}
\usepackage{hyperref}
\usepackage{amsmath,amssymb}
\usepackage{parskip}
\usepackage{titlesec}

\titleformat{\section}{\large\bfseries}{}{0em}{}
\titleformat{\subsection}{\normalsize\bfseries}{}{0em}{}
\titlespacing{\section}{0pt}{6pt}{2pt}
\titlespacing{\subsection}{0pt}{4pt}{1pt}
\hypersetup{colorlinks=true, linkcolor=blue, urlcolor=blue, citecolor=blue}
\setlength{\parskip}{3pt}

\begin{document}

\begin{center}
  {\Large \textbf{CS590: Socially Cognizant Robotics}} \\[5pt]
  {\Large \textbf{Deliverable 2}} \\[12pt]
  \textbf{Group Members:} \\[3pt]
  Rohit Bellam, rsb204 \\
  Karim Smires, ks1686 \\[8pt]
  March 4th, 2026
\end{center}

\vspace{2pt}

\section{1. Overview}

Following the framework of Kress-Gazit et al.\ in \textit{Robot
  Learning as an Empirical Science: Best Practices for Policy
Evaluation}~\cite{bestpractices}, we analyze the evaluation
methodology of \textit{Dream to Control} (Hafner et al.,
ICLR 2020). We identify missing metrics, insufficient scene
diversity, and unreported failure modes that would strengthen the
understanding of Dreamer's capabilities and limitations.

\section{2. Additional Metrics}

\subsection{2.1 Lack of Success Criteria and Sub-Goal Decomposition}
Dreamer reports performance mainly through \textbf{episode return}
averaged over 5 seeds---analogous to the ``success rate only''
problem identified by Kress-Gazit et al. No binary
success/failure threshold is defined for any task, making it
impossible to determine what fraction of episodes produce stable
locomotion versus partial failures. Moreover, Dreamer's tasks are
multi-phase (e.g., Hopper Hop requires standing, balancing, then
hopping), but no sub-goal rubric is reported. As the best
practices paper highlights with its Pancake example (Table 1),
sub-goal tracking can reveal that a policy with low overall success
may still reliably complete early stages, providing a
gradient for understanding partial progress.

\subsection{2.2 Statistical Analysis}
While Dreamer reports mean and standard deviation over 5 seeds, no
formal statistical comparison is performed between methods. As
Kress-Gazit et al.\ emphasize, point estimates can be
misleading---bootstrap confidence intervals or stratified
permutation tests would quantify whether Dreamer's advantage over
PlaNet or D4PG is statistically significant on individual tasks,
rather than relying on visual inspection of overlapping error bars.
The number of evaluation episodes per checkpoint is also unreported,
a concern raised by the best practices paper.

\subsection{2.3 Trajectory Smoothness and Formal Specifications}
For real-world deployment, \textit{how} a robot moves matters as much
as whether it succeeds. The best practices paper recommends SPARC
(SPectral ARC length) to quantify trajectory smoothness---relevant for Dreamer's
locomotion tasks, where jerky or oscillating trajectories may achieve
high returns in simulation but would fail on physical hardware.
Additionally, Signal Temporal Logic (STL) robustness could
automatically evaluate safety constraints (e.g., joint torque limits,
ground contact forces), quantifying \textit{how close} the agent is to
violating a specification.

\section{3. Additional Scenes and Environments}

\subsection{3.1 Limited Visual Diversity and No Out-of-Distribution Testing}
All 20 DeepMind Control Suite tasks use simple, uniform backgrounds
with consistent lighting at $64 \times 64$ resolution. There is no
evaluation of robustness to visual perturbations such as lighting changes,
camera shifts, or background clutter---conditions that benchmarks
like Distracting Control Suite were designed to test.
Since Dreamer's world model is trained via pixel reconstruction,
visual simplicity likely inflates performance by making reconstruction
trivially easy; it remains unclear whether the latent dynamics would
remain accurate under richer visual conditions.

\subsection{3.2 Narrow Task Diversity and No Sim-to-Real Transfer}
The 20 DMC tasks are predominantly locomotion-based. There are no
manipulation tasks involving object interaction, no navigation with
obstacles, and no multi-agent scenarios. Analogous to the ``forest
vs.\ bug trap'' distinction in the milestone description, the DMC
tasks---while varied in dynamics---do not test for qualitatively
different challenges such as deceptive reward landscapes that induce
local optima in
the learned policy. Furthermore, despite Dreamer's emphasis on
computational efficiency for robotics, no sim-to-real experiments
are conducted. Even a simple domain transfer test would better
support the paper's claims of practical relevance.

\section{4. Failure Modes}

\subsection{4.1 Reported Failure Modes}
The paper acknowledges two failure modes through ablations: (1)
\textbf{shortsighted behavior} when the value model is removed,
causing poor performance on long-horizon tasks like Acrobot Swingup,
and (2) \textbf{representation insufficiency} when using reward-only
prediction, which fails to learn useful latent dynamics. The
Atari/DMLab experiments further suggest the reconstruction-based
world model struggles with visually complex environments.

\subsection{4.2 Unreported Failure Modes}
Several important failure modes are neither analyzed nor discussed:

\textbf{World model compounding error:} As the imagination horizon
increases, transition model prediction errors accumulate. While
Figure 4 shows robustness to horizon length in aggregate, there is no
analysis of when or how predictions degrade, or how compounding error
affects specific tasks.

\textbf{Sensitivity to initial conditions:} The best practices paper
demonstrates (Examples 1--2) that learned policies are highly
sensitive to initial conditions. Dreamer reports no
per-initial-condition analysis despite using randomized initial
states, making it unclear whether strong average performance masks
consistent failure from certain configurations.

\textbf{Exploration failures:} Dreamer uses simple Gaussian noise
($\sigma = 0.3$) for exploration. In sparse-reward tasks (Cartpole
Swingup Sparse, Acrobot Swingup), this may be insufficient to
discover rewarding trajectories, yet no analysis of exploration
adequacy or comparison to structured exploration strategies is
provided.

\textbf{Observation noise and partial observability:} Although
formulated as a POMDP, all evaluations use clean, noise-free rendered
images. Robustness to sensor noise, dropped frames, or occluded
observations---conditions ubiquitous in real robotic
systems---is untested.

\textbf{Distribution shift (objective mismatch):} As the policy
improves, it visits state regions underrepresented in the replay
buffer. The world model may then make poor predictions precisely
where accurate imagination matters most, potentially causing policy
degradation---a known issue in model-based RL that is not analyzed.

\vspace{2pt}
\begin{thebibliography}{9}
  \small
  \bibitem{bestpractices} H.\ Kress-Gazit et al., ``Robot Learning as
  an Empirical Science: Best Practices for Policy Evaluation,''
  \textit{arXiv:2409.09491}, 2024.
\end{thebibliography}

\end{document}
